\subsection{Computational verification: exact coordinates, certificates, and chronology}
\label{subsec:computational-verification}

The computational literature does not use one unnamed Collatz map.  In this
subsection
\[
 U(n)=
 \begin{cases}
  3n+1,&n\text{ odd},\\
  n/2,&n\text{ even},
 \end{cases}
 \qquad
 T(n)=
 \begin{cases}
  (3n+1)/2,&n\text{ odd},\\
  n/2,&n\text{ even}
 \end{cases}
\]
denote respectively the unshortened and shortened maps.  One odd $T$-step is
two $U$-steps and one even $T$-step is one $U$-step.  Barina also uses the
auxiliary chart
\[
 T_1(n)=
 \begin{cases}
  (n+1)/2,&n\text{ odd},\\
  3n/2,&n\text{ even}.
 \end{cases}
\]
Its exact branchwise bridges are
\[
 T(n)=
 \begin{cases}T_1(n+1)-1,&n\text{ odd},\\ n/2,&n\text{ even},\end{cases}
 \quad
 T_1(n)=
 \begin{cases}T(n-1)+1,&n\text{ odd},\\3n/2,&n\text{ even}.
 \end{cases}
\]
They move the additive term between coordinate charts; they do not identify
the maps or their clocks.  In particular, the 2025 paper's page-4
\emph{trajectory on $n+1$}, $(7,4,6,9,5,3,2)$, is the shifted-coordinate
presentation associated with the trajectory from $12$, not the $T$-orbit
of $13$ \cite[pp.~2, 4]{Barina2025}.

\begin{proposition}[Barina's block map, with fibres]
\label{prop:Barina-block-map}
For $x\in\Npos$, put
\[
 \alpha=\nu_2(x+1),\qquad u=\frac{x+1}{2^\alpha},\qquad
 y=3^\alpha u-1,\qquad \beta=\nu_2(y),
\]
and $B(x)=y/2^\beta$.  Then $B(x)$ is a positive odd integer and
\begin{equation}
\stepcounter{equation}
\label{eq:Barina-block}
 B(x)=T^{\alpha+\beta}(x).
\tag{CV.1}
\end{equation}
For every odd $z>0$, the fibre $B^{-1}(z)$ is exactly the set of
\begin{equation}
\stepcounter{equation}
\label{eq:Barina-fibre}
 x=2^\alpha\frac{2^\beta z+1}{3^\alpha}-1
\tag{CV.2}
\end{equation}
over integers $\alpha\geq0,\beta\geq1$ satisfying
$3^\alpha\mid2^\beta z+1$.  Hence $B$ is surjective, has infinite fibres,
and has no left or two-sided inverse after the block boundary and elapsed
clock have been discarded.
\end{proposition}

\begin{proof}
The integer $u$ is odd.  For $0\leq j\leq\alpha$, direct induction gives
\[
 T^j(x)=3^j\frac{x+1}{2^j}-1.
\]
The state at $j=\alpha$ is $y$; the following $\beta$ steps divide it by
two, proving \eqref{eq:Barina-block}.  Conversely,
\eqref{eq:Barina-fibre} gives $x+1=2^\alpha u$, with
$u=(2^\beta z+1)/3^\alpha$ odd, and
$3^\alpha u-1=2^\beta z$.  Thus the two valuations are exactly
$\alpha,\beta$ and $B(x)=z$.  The section $z\mapsto2z$ proves
surjectivity.  Taking $\alpha=0$ in \eqref{eq:Barina-fibre} gives the
distinct points $2^\beta z$, $\beta\geq1$, in the same fibre.
\end{proof}

On the odd $2$-adic integers with normalized Haar measure,
$\Pr(\alpha=a)=2^{-a}$ and
$\Pr(\beta=b\mid\alpha=a)=2^{-b}$ for $a,b\geq1$.  Multiplication by the
odd unit $3^a$ proves the conditional statement, so the two valuations are
independent and
\[
 \mathbb E(\alpha+\beta)=4.
\]
This is the exact measure-theoretic content of the articles' phrase
\emph{average four iterations}; it is not an assertion that every finite
interval has raw arithmetic mean four \cite[pp.~3--5]{Barina2021}.

\begin{sourceissue}[the exceptional start $1$ in the printed algorithms]
\label{sourceissue:Barina-one-endpoints}
Equation~\eqref{eq:Barina-block} has $B(1)=1$, representing the shortened
cycle $1\mapsto2\mapsto1$.  The printed endpoints differ by version.  The
2021 Algorithm~1 uses strict \texttt{until n<n0} and loops at $n_0=1$.
Its Algorithm~2 uses \texttt{until n=1}: it terminates after one complete
$1\mapsto1$ macro, but a literal delay accumulator gives $2$ rather than
the conventional stopping time $0$.  Both 2025 Algorithms~1 and~2 use the
strict endpoint and loop at $1$ \cite[pp.~4--5]{Barina2021}
\cite[pp.~4--5]{Barina2025}.  The exact repairs are a separate base case, or
the domain $n_0>1$ where only descent is required.  The production inner
loop tests starts $3\bmod4$, so this exceptional input is not encountered
there.
\end{sourceissue}

\begin{proposition}[the dyadic affine-prefix coordinate]
\label{prop:computational-dyadic-prefix}
Fix $k\geq0$, $h\geq0$, and $0\leq r<2^k$, and let
\[
 a_j(r)=\#\{0\leq i<j:T^i(r)\text{ is odd}\}.
\]
For every $0\leq j\leq k$,
\begin{equation}
\stepcounter{equation}
\label{eq:computational-dyadic-prefix}
 T^j(2^kh+r)=3^{a_j(r)}2^{k-j}h+T^j(r).
\tag{CV.3}
\end{equation}
In particular,
\begin{equation}
\stepcounter{equation}
\label{eq:computational-dyadic-endpoint}
 T^k(2^kh+r)=3^{a_k(r)}h+T^k(r).
\tag{CV.4}
\end{equation}
\end{proposition}

\begin{proof}
For $j<k$, the high term in \eqref{eq:computational-dyadic-prefix} is even.
The full state and the low state $T^j(r)$ therefore have the same parity and
take the same next branch.  Applying that branch proves the induction.  This
is simultaneously Oliveira e Silva's Proposition~1 and Barina's equation~(6),
after their maps and variables have been typed rather than identified
\cite[p.~191]{OliveiraSilva2010}\cite[p.~5]{Barina2021}.
\end{proof}

Write $a=a_k(r)$, $d=T^k(r)$, $L=3^a-2^k$, and $R=r-d$.  The exact
necessary-and-sufficient condition for
\[
 T^k(2^kh+r)<2^kh+r\qquad(h=1,2,\ldots)
\]
is
\begin{equation}
\stepcounter{equation}
\label{eq:computational-uniform-descent}
 (L<0\text{ and }L<R)\quad\text{or}\quad(L=0\text{ and }R>0).
\tag{CV.5}
\end{equation}
Indeed the inequality is $Lh<R$, whose most restrictive positive $h$ is
$h=1$ when $L<0$; it is constant when $L=0$, and eventually fails when
$L>0$.  The generator's \texttt{is\_convergent} implements only the first
disjunct.  It is conservative in the equality case; it does not falsely
certify it.

If $r'<r$ and
\[
 a_k(r')=a_k(r),\qquad T^k(r')=T^k(r),
\]
then \eqref{eq:computational-dyadic-endpoint} proves, for every $h$,
\begin{equation}
\stepcounter{equation}
\label{eq:computational-endpoint-coalescence}
 T^k(2^kh+r')=T^k(2^kh+r),\qquad2^kh+r'<2^kh+r.
\tag{CV.6}
\end{equation}
Thus a dead sieve bit has a finite descent, coalescence, inherited-prefix, or
restricted-predecessor certificate.  A live bit is merely the complement of
the implemented certificates.  It is not a residue class proved not to
converge; the 2025 page-6 wording \emph{do not converge} must be read at this
finite algorithmic scope \cite{Barina2025}.

\begin{sourceissue}[cardinality and generator-history defects]
\label{sourceissue:Barina-sieve-cardinality}
Within one half-open work unit of width $2^{40}$, a fixed class modulo
$2^k$ contains $2^{40-k}$ starts, while there are $2^k$ classes.  The 2025
Section~3.4 reverses these quantities twice.  The 2021 page-6 example and
\path{src/worker/worker.c:336--369} use the correct $2^{40-34}$.

The separate \texttt{collatz-sieve} history has a dated implementation
  defect.  Commit
  \nolinkurl{49c5524d8d42867c866e863030a8e5d86237bf17}
(10 April 2020) passed the odd-step count \texttt{c\_} as both the third and
fourth arguments of \texttt{is\_convergent}; the fourth argument had to be
the endpoint \texttt{d\_}.  Commit
  \nolinkurl{b3d2f61d89eec55008d61c32577a9e8c8b3874ab}
(6 April 2023) repaired it.  The baseline committed maps predate the faulty
interval, while the \texttt{hesieve} and \texttt{h2esieve} maps postdate the
repair.  The defect governs regeneration from intermediate source history;
it is not silently attributed to blobs outside that interval.
\end{sourceissue}

The first exact lower-start reductions are
\[
 T(2q+1)=3q+2,
\]
and
\begin{equation}
\label{eq:computational-four-transports}
 T^3(8q+3)=9q+4,\qquad
 T^6(64q+7)=81q+10,\qquad
 T^7(128q+95)=243q+182.
\end{equation}
Their odd-step counts are respectively $1,2,4,5$.  Every source on the left
is strictly smaller than its target.  The first three families omit $37$ of
$81$ initial residues: the first two give
$\{2,4,5,8\}\bmod9$, their lifts modulo $27$ are
\[
 \{2,4,5,8,11,13,14,17,20,22,23,26\},
\]
and the $81q+10$ class is disjoint from them.  The fourth target is already
$2\bmod3$ as an initial residue, but is not redundant when it occurs as an
intermediate low-coordinate state.

\begin{proposition}[the Hercher coalescence lift]
\label{prop:Hercher-computational-coalescence}
Let $x$ be odd and put
\[
 \alpha=\nu_2(x+1),\quad u=(x+1)/2^\alpha,
 \quad y=3^\alpha u-1,\quad\beta=\nu_2(y).
\]
If $\beta\geq2$ and $m=(x-1)/2$, then
\begin{equation}
\stepcounter{equation}
\label{eq:Hercher-local-coalescence}
 T^{\alpha+\beta-1}(m)=T^{\alpha+\beta}(x)=y/2^\beta.
\tag{CV.7}
\end{equation}
If $x=T^s(r)$, the prefix contains $c$ odd steps, and $m<r$, then for every
$N=2^kh+r$ the lower lifted start
\[
 Q=3^c2^{k-s-1}h+m
\]
satisfies $Q<N$ and its trajectory coalesces with that of $N$.
\end{proposition}

\begin{proof}
Since $y\equiv0\pmod4$, direct iteration of $m$ gives
\eqref{eq:Hercher-local-coalescence}.  The affine prefix has nonnegative
intercept, positive when $c>0$.  From $m<r$ one has $x<2r$, hence
$3^c/2^s<2$.  This is exactly the strict high-coefficient comparison for
$Q<N$; \eqref{eq:computational-dyadic-prefix} and
\eqref{eq:Hercher-local-coalescence} then prove coalescence.
\end{proof}

The additional Oliveira--Roosendaal example is likewise an identity:
\begin{equation}
\stepcounter{equation}
\label{eq:Oliveira-Roosendaal-coalescence}
 T^6(64q+15)=T^6(64q+14)=T^5(32q+7)=81q+20.
\tag{CV.8}
\end{equation}

\begin{theorem}[finite closure of the recursive high-coefficient check]
\label{thm:Barina-recursive-k34}
Consider \texttt{collatz-sieve}'s restricted recursive closure generated by
the four inverse transports above.  For every invocation arising in its
packed implementation domain $k\leq34$, a returned lower representative
$P<r$ has a uniformly lower full-width lift.  This assertion is finite in
$k$; it does not prove an unbounded recursive invariant.
\end{theorem}

\begin{proof}
A recursive word with $d$ odd steps and $t$ total shortened steps gives
\[
 T^t(P)=L.
\]
For the four primitive words, $(d,t,H)$ is
\[
 (1,1,1),\ (2,3,5),\ (4,6,73),\ (5,7,211),
\]
and the inverse is the exact floor map
\[
 \phi(L)=\frac{2^tL-H}{3^d}
         =\left\lfloor\frac{2^tL}{3^d}\right\rfloor.
\]
A nonempty composite word has
\begin{equation}
\stepcounter{equation}
\label{eq:recursive-rounded-word}
 P=\frac{2^tL-H_w}{3^d}=\frac{2^t}{3^d}L-h_w,
 \qquad0<h_w<\ell(w)\leq d.
\tag{CV.9}
\end{equation}

Suppose $L=T^s(r)$ after $c$ odd steps.  Substitution in
\eqref{eq:computational-dyadic-prefix} shows that the unique
coefficient-correct full lift is
\[
 Q=3^{c-d}2^{k-s+t}h+P,
\]
and $Q<2^kh+r$ for every $h>0$ precisely when $P<r$ and
\begin{equation}
\stepcounter{equation}
\label{eq:recursive-current-coefficient}
 2^t3^{c-d}\leq2^s.
\tag{CV.10}
\end{equation}
For the recursive branch applied after
Proposition~\ref{prop:Hercher-computational-coalescence}, the corresponding
condition, at the preceding prefix $(s_0,c_0)$, is
\begin{equation}
\stepcounter{equation}
\label{eq:recursive-coalescence-coefficient}
 2^t3^{c_0-d}\leq2^{s_0+1}.
\tag{CV.11}
\end{equation}
The C function returns only $P$, retains neither $d$ nor $t$, and tests
neither inequality.

An exact finite reduction supplies the missing check.  In the relevant
exponent boxes the smallest coefficient strictly above one is
\begin{equation}
\stepcounter{equation}
\label{eq:recursive-minimum-unsafe}
 C_{\min}=\frac{3^{12}}{2^{19}}
 =\frac{531441}{524288},
 \qquad C_{\min}-1=\frac{7153}{524288}.
\tag{CV.12}
\end{equation}
If \eqref{eq:recursive-current-coefficient} failed while $P<r$,
equation~\eqref{eq:recursive-rounded-word} and the affine prefix imply
$P(C-1)<h_w<34$ and $r-P\leq33$; hence $P\leq2492$.  For
\eqref{eq:recursive-coalescence-coefficient}, the extra half-step contributes
less than $1/2$, so an unsafe witness has $P\leq2528$ and $r-P\leq34$.

The deterministic certificate exhausts these proved supersets.  In the
current-state branch it checks $7105$ restricted word states, $1557$
identity endpoint pairs, $2983$ nonidentity common-endpoint pairs, and all
$2585$ budget-admissible pairs.  In the coalescence branch it checks $7216$
word states, $1003$ common-endpoint pairs, including $480$ identity and $523$
nonidentity pairs, and all $487$ budget-admissible nonidentity pairs.  Every
required comparison satisfies \eqref{eq:recursive-current-coefficient} or
\eqref{eq:recursive-coalescence-coefficient}.  The finite minimum
\eqref{eq:recursive-minimum-unsafe}, rather than an asymptotic claim, is what
closes the implementation domain.
\end{proof}

\begin{sourceissue}[remaining code and witness boundaries]
\label{sourceissue:Barina-code-boundaries}
The name \texttt{smallest\_predecessor} means the minimum only in the
four-map recursive closure, not among all Collatz predecessors.  At residue
zero, the unsigned expression $(L_{\rm last}-1)/2$ wraps; it does not clear
zero but has no positive-integer predecessor meaning.  The effective
$k\leq34$ domain is the packed-field assertion $b<2^{34}$, not a function
precondition; disabling assertions outside it permits the 34-bit residue
field to alias.  The \texttt{ctzu64} helper assumes that
\texttt{unsigned long} has 64 bits and traps on LLP64.  Finally, a dead-map
bit stores no predecessor word, step counts, or elimination-reason witness,
so the bit alone cannot reconstruct its v3 certificate.
\end{sourceissue}

\begin{proposition}[finite convergence from certified descent]
\label{prop:finite-descent-induction}
If, for every $2\leq n\leq X$, a computation exhibits a finite iterate
$T^{j(n)}(n)<n$, then every positive start through $X$ reaches $1$.
\end{proposition}

\begin{proof}
Strong induction on $n$.  The displayed iterate is a smaller positive start,
so its orbit reaches $1$ by the induction hypothesis; concatenate the two
finite orbit segments.  The base $n=1$ is immediate.
\end{proof}

This proposition explains why sieve elimination and stopping-time descent
can establish a finite range without individually iterating every start to
one.  It does not establish that a reported machine run occurred.

\begin{center}
\begin{tabular}{@{}lll@{}}
\toprule
Source & exact reported range & status in this edition\\
\midrule
Oliveira e Silva (2010) & $n\leq20\cdot2^{58}$ & sourced finite result\\
Barina (2021) & $n<2^{68}$ & sourced finite result\\
Barina (2025) & $n<2^{71}$, 15 Jan.\ 2025 & sourced finite result\\
Project API, 28 Aug.\ 2026 & $2075\cdot2^{60}$ & later dynamic assertion\\
\bottomrule
\end{tabular}
\end{center}

Oliveira e Silva's printed inclusive endpoint is
\[
 20\cdot2^{58}=5764607523034234880.
\]
The twenty directly processed intervals are naturally half-open, but the
boundary itself is even and maps to the already covered $10\cdot2^{58}$;
hence the inclusive theorem follows.  The chapter reports a depth-46 pruned
parity tree with $269949796982$ surviving leaves, $2^{12}$-start blocks,
variable-step tables with $\Delta=14$ and tolerance $\tau=10$, two
independent workers per item, and exact agreement of candidate data and a
cyclic-redundancy checksum.  One mismatch in about 81 CPU-years was
invalidated and recomputed \cite[pp.~189--198]{OliveiraSilva2010}.  The
underlying worker source, completed-work ledger, and duplicate outputs were
not acquired, so the range is not relabelled an F031 reproduction.

The 2021 article reports the $2^{68}$ run from September 2019 to May 2020
\cite[p.~7]{Barina2021}.  The 2025 Table~10 records the strict boundary
$n<2^{71}$ and completion date; its surrounding phrase \emph{up to} does not
change the half-open work units
\[
 [j2^{40},(j+1)2^{40}).
\]
The boundary $2^{71}$ itself is a power of two and independently descends,
but was not in the first $2^{31}$ units.  Reported CPU/GPU years, throughput,
and speedups are measurements on named configurations, not
hardware-independent complexity theorems.  The page-2 throughput convention
also counts starts eliminated by sieve certificates whose trajectories were
not individually run to one \cite{Barina2025}.

The dated API snapshot reports
\[
 2075\cdot2^{60}=2392312122059207475200
 \approx2^{71.0188956211},
\]
with the first incomplete unit beginning at
\[
 2175975546\cdot2^{40}=2392510414583230365696.
\]
It is later than the paper, is not a completed-work ledger, and is not
backdated into the 2025 theorem \cite{BarinaStatus2026}.

\begin{sourceissue}[the public-code validation boundary]
\label{sourceissue:Barina-public-code-validation}
The 2025 paper pins the repository \nolinkurl{xbarin02/collatz} at commit
\linebreak[3]\nolinkurl{53c2a0608075d6fe3f10cc6eeeaf50e400c86338}.
It does not pin the
separate \texttt{collatz-sieve} dependency.  The bootstrap performs unpinned
clone-or-pull operations.  The edition's sieve commit and map hashes are
therefore separately attributed reconstruction pins, not historical facts
invented for the run \cite{BarinaCode2025}.

The default GPU kernel is the Git symlink \texttt{kernel.cl} to
\texttt{kernel32.cl}.  A GPU overflow zeroes a checksum lane; the host aborts
the unit and mclient attempts a CPU rerun.  GMP continuation exists only when
that CPU worker was compiled with \texttt{\_USE\_GMP}.  The generic bootstrap
does so, whereas the Lumi and Meta GPU-only scripts build no CPU worker.
Thus GPU--CPU--GMP is a conditional deployment chain, not GPU
multiprecision or a universal local fallback.

The optional/test \texttt{kernel32-precalc.cl} has a sieve-dependent
\texttt{continue} immediately before a work-group barrier.  A multi-lane
group with differing sieve bits violates barrier convergence.  This defect
is scoped to that variant and is not transferred to the default kernel.

The server first rejects stale returns with the wrong stored client ID and
rejects zero or non-whitelisted checksum prefixes.  For an accepted current
return, however, it marks the unit complete before comparing a previous
checksum, then logs and overwrites a mismatch; maximum offsets are handled in
the same order.  Mclient transmits the $\alpha$-sum and discards the
$\beta$-sum.  The verification utility recomputes only the trajectory maximum
at the stored offset, not the work unit.  These are consistency signals, not
Oliveira e Silva's independent duplicate-run protocol or a cryptographic
certificate.

No acquired immutable ledger binds every unit to its code tree, map hash,
binary, compiler, OpenCL environment, and return value.  Selected regression
tests and public source permit code-level audit; they do not independently
authenticate the full historical $2^{71}$ computation.
\end{sourceissue}

\begin{proposition}[shortened and unshortened maximum excursions]
\label{prop:shortened-unshortened-maxima}
For an odd $n>1$, if the orbit maxima are finite, then
\begin{equation}
\stepcounter{equation}
\label{eq:shortened-unshortened-maxima}
 M_U(n)=2M_T(n).
\tag{CV.13}
\end{equation}
The nontrivial record-start sequence is the same under both conventions,
although its maximum values differ by the factor two.
\end{proposition}

\begin{proof}
An odd shortened state $x>1$ cannot be the shortened maximum because
$(3x+1)/2>x$.  A maximal shortened state cannot have an even predecessor,
which would be twice as large.  It is therefore the output $y=(3x+1)/2$ of
an odd predecessor, and the unshortened orbit contains the intermediate
state $2y$.  Every state omitted by shortening is twice a shortened state,
proving \eqref{eq:shortened-unshortened-maxima}.  No even $n>2$ begins a new
unshortened record, because the lower odd start $n-1$ has first image
$3n-2>n$.  Thus record starts are odd and multiplication of every maximum by
two preserves their order.
\end{proof}

The 2025 abstract says four new records, while Section~6 lists five and calls
them new.  The first listed start, $274133054632352106267$, already appears
in the 2021 paper.  The exact reconciliation is four records new after 2021
and five cumulative project records; Section~6's phrase \emph{five new} is
imprecise rather than an unexplained numerical contradiction
\cite[p.~7]{Barina2021}\cite[pp.~12--13]{Barina2025}.  Finite record data are
consistent with the Lagarias--Weiss prediction but do not prove its
asymptotic form.

\begin{proposition}[Roosendaal's modular maximum theorem]
\label{prop:Roosendaal-mod36}
Let $N>2$ be odd.  If its unshortened maximum $q=M_U(N)$ is finite and
$q>3N+1$, then
\[
 q\equiv16\pmod{36}.
\]
\end{proposition}

\begin{proof}
The maximum $q$ is even and cannot be $2\bmod4$, since then the next two
states $q/2$ and $3q/2+1$ exceed it.  Because $q>3N+1$, it is neither the
start nor its first odd output; halving cannot create a maximum, so
$q=3p+1$ for an odd predecessor and $q\equiv1\bmod3$.  The remaining classes
modulo $36$ are $4,16,28$.  If $q=36a+4$, its forced preceding states include
$48a+4>q$.  If $q=36a+28$, they include $48a+36>q$.  Only $16\bmod36$
remains \cite{Roosendaal2026}.
\end{proof}

The page's following sentence applies this congruence to $P_i$, but the
theorem concerns $M_U(P_i)$, not the starting values.  Its Conjecture~3 also
prints $N<1$ against the surrounding $N>1$ domain.  Neither defect is
silently normalized.  Theorems~2--3 each begin with a residue-bound
hypothesis that is not proved on the page; this edition does not establish
that hypothesis and therefore does not import either conclusion.

Oliveira e Silva additionally studies the one-branch-per-iterate map
\begin{equation}
\stepcounter{equation}
\label{eq:Oliveira-generalized-five}
 C(n)=
 \begin{cases}
  n/2,&n\equiv0\pmod2,\\
  n/3,&n\not\equiv0\pmod2,\ n\equiv0\pmod3,\\
  5n+1,&\gcd(n,6)=1.
 \end{cases}
\tag{CV.14}
\end{equation}
Its clock is not the $T$- or $U$-clock.

\begin{sourceissue}[Oliveira e Silva's Proposition~2 is false as printed]
\label{sourceissue:Oliveira-proposition-two}
The proposition asks the first $i$ iterates of $m_i$ to determine counts
$k_2,k_3,k_5$, puts
\[
 r_i=2^{k_2}3^{k_3},\qquad
 m_i=n_0\bmod r_i,\qquad n_i=[n_0/r_i],
\]
and asserts
\[
 C^i(n_0)=5^{k_5}n_i+C^i(m_i).
\]
At $i=1,n_0=1$, $r_i=1$ forces $m_i=0$, whose division-by-two branch demands
$r_i=2$; $r_i=2$ or $3$ forces $m_i=1$, whose $5n+1$ branch demands
$r_i=1$.  Thus the circular definition has no fixed point.  If one instead
takes the branch of $n_0$ first, then $r_i=1,m_i=0,n_i=1$ and the asserted
equation reads $6=5$.  No domain qualification repairs it
\cite[p.~200]{OliveiraSilva2010}.

The chapter defines $[x]$ as the largest integer not smaller than $x$, which
is ceiling language.  Proposition~1's unique dyadic decomposition requires
$\lfloor n_0/2^i\rfloor$.  Replacing the bracket by a floor repairs that use;
it does not repair Proposition~2.
\end{sourceissue}

\begin{proposition}[the corrected $6^i$ branch-cylinder coordinate]
\label{prop:Oliveira-branch-cylinder}
For $i\geq0$, put
\[
 a_i=n_0\bmod6^i,\qquad
 Q_i=\left\lfloor\frac{n_0}{6^i}\right\rfloor.
\]
Let $k_2,k_3,k_5$ be the branch counts in the first $i$ iterates of $a_i$,
set $R_i=2^{k_2}3^{k_3}$, and
\[
 N_i=\frac{n_0-a_i}{R_i}=\frac{6^i}{R_i}Q_i.
\]
Then $N_i$ is an integer, $n_0$ and $a_i$ have the same first $i$ branches,
and
\begin{equation}
\stepcounter{equation}
\label{eq:Oliveira-branch-cylinder}
 C^i(n_0)=5^{k_5}N_i+C^i(a_i).
\tag{CV.15}
\end{equation}
The representative $a_i$ is a branch-cylinder representative, generally not
the canonical remainder modulo $R_i$.
\end{proposition}

\begin{proof}
If $x\equiv y\pmod{6^j}$, the two points take the same branch.  After
division by two their difference is divisible by $2^{j-1}3^j$, after
division by three by $2^j3^{j-1}$, and after $5x+1$ by $6^j$; every image
pair is congruent modulo $6^{j-1}$.  Induction gives the common length-$i$
branch word.

For a fixed word $w$, induction on its branches gives an integer $B_w$ with
\[
 C^i(x)=\frac{5^{k_5}x+B_w}{R_i}
\]
throughout its cylinder.  A division branch multiplies $R_i$ by two or
three, while a $5x+1$ branch replaces $B_w$ by $5B_w+R_i$.  Since
$R_i\mid6^i$, subtracting the formula at $n_0$ and $a_i$ proves
\eqref{eq:Oliveira-branch-cylinder}.  Reducing $a_i$ modulo $R_i$ is valid
only after a separate proof that the reduced point retains the same branch
word.
\end{proof}

The source reports exhaustive computations through $10^{16}$ for
\eqref{eq:Oliveira-generalized-five} and $10^{14}$ for the analogous
$C_{7,\{2,3,5\}}$, taking about five and two CPU-months respectively on one
450 MHz Pentium~III \cite[pp.~200--204]{OliveiraSilva2010}.  It does not say
that the replicated-worker protocol of the $3x+1$ run was used here.  Its
finite pruning counts, affine tables, transition matrix, and stationary
vector are reproducible finite calculations.  Section~4 expressly calls the
passage from a Chernoff upper bound to an asymptotic equality non-rigorous;
the printed values
\[
 \rho_5\approx1.442762198279,\qquad
 \rho_7\approx2.215508001593
\]
remain conjectured record-growth exponents, not proved rates.

\begin{nonclaim}
The deterministic F031 certificate pins the exact manifestations, code
trees, bundles, selected maps, and route artifacts; reproduces the basic
$2^{16}$ sieve byte for byte; proves
Theorem~\ref{thm:Barina-recursive-k34}; and checks the map, fibre, affine,
transport, maximum, modular, and corrected branch-cylinder identities above.
It does not reproduce a historical distributed run, authenticate a
completed-work database, independently regenerate the full $2^{24}$ or
$2^{34}$ maps, certify hardware or compilers, justify the Markov asymptotics,
or prove universal convergence. The eight legacy unresolved-dependency records
are audit-schema records rather than mathematical claims; this section does not
use them to preserve or conditionalize an unproved conclusion.
\end{nonclaim}
